% !TEX root = ../../../main.tex
% 来源：中学数学实验教材·第二册（几何）
% 章节：直线形 / 特殊四边形
% 原始文件：3.tex，行 3851-3884
% 类型：tikz
% ---
\begin{tikzpicture}
\tkzDefPoints{0/0/B, 3/0/C, 3/2/D, 0/2/A}
\tkzDrawPolygon[thick](A,B,C,D)
\tkzDrawLines[add=0 and .5, dashed](B,A A,D D,C C,B)
\tkzDefPointWith[linear, K=1.5](B,A)  \tkzGetPoint{A'}
\tkzDefPointWith[linear, K=1.5](A,D)  \tkzGetPoint{D'}
\tkzDefPointWith[linear, K=1.5](D,C)  \tkzGetPoint{C'}
\tkzDefPointWith[linear, K=1.5](C,B)  \tkzGetPoint{B'}

\tkzDefLine[bisector](D,A,A') \tkzGetPoint{A''}
\tkzDefLine[bisector](C,D,D') \tkzGetPoint{D''}
\tkzDefLine[bisector](B,C,C') \tkzGetPoint{C''}
\tkzDefLine[bisector](A,B,B') \tkzGetPoint{B''}
\tkzInterLL(A,A'')(D,D'')  \tkzGetPoint{H}
\tkzInterLL(D,D'')(C,C'')  \tkzGetPoint{G}
\tkzInterLL(B,B'')(C,C'')  \tkzGetPoint{F}
\tkzInterLL(A,A'')(B,B'')  \tkzGetPoint{E}
\tkzDrawPolygon[thick](E,F,G,H)
\tkzInterLL(A,E)(C,B)  \tkzGetPoint{P}
\tkzDrawSegments[dashed](E,P P,B)

\tkzLabelPoints[above](H,D)
\tkzLabelPoints[below](P,F,B)
\tkzLabelPoints[left](E,A)
\tkzLabelPoints[right](G,C)

\tkzMarkAngles[mark=none, size=.4](E,A,B A,B,E D,A,H B,C,F F,B,C H,D,A B,P,E)
\tkzLabelAngle[pos=.6](E,A,B){1}
\tkzLabelAngle[pos=.6](A,B,E){2}
\tkzLabelAngle[pos=.6](D,A,H){3}
\tkzLabelAngle[pos=.6](B,C,F){4}
\tkzLabelAngle[pos=.6](F,B,C){5}
\tkzLabelAngle[pos=.6](H,D,A){6}
\end{tikzpicture} 
